% End-to-end comparison with Llumnix on A100 (Astrolabe as default)
\subsection{End-to-end Comparison with Llumnix}\label{section:llumnix_comparison}

We compare Astrolabe with full Llumnix~\cite{sun2024llumnix}, which performs dynamic load rebalancing via live KV-cache migration. Our A30 cluster (single-GPU-per-host, 25\,Gbps Ethernet) is unsuitable for migration due to network constraints. We therefore use a 2-node cluster with 4$\times$A100-40GB GPUs per node, analogous to Llumnix's original 4-node A10 testbed. Llumnix moves KV-cache via Ray RPC by default and does not support chunked prefill (CP) in the evaluated configuration. To keep the comparison faithful to the released system, we keep Llumnix's default Ray-RPC migration path and do not retune its migration policy for our workload; we change only deployment parameters needed to run the same model and instance count on the A100 cluster. To isolate the scheduler's contribution from chunked-prefill backend effects, we evaluate Astrolabe with and without CP against Llumnix across an identical QPS sweep. With larger A100 memory (40\,GB vs 24\,GB), we raise the batch size to 96. All configurations serve Llama-2-7B across 8 instances (4 per node, TP=1 for each), using FlashInfer as backend.

\begin{figure}[!tb]
    \centering
    \includegraphics[width=0.90\linewidth]{figure/revision/llumnix_comparison_v2_po2.png}
    \caption{Astrolabe with and without chunked prefill vs full Llumnix on the A100 cluster}
    \Description{Throughput, mean request latency, mean token latency, and scheduling overhead versus QPS for Astrolabe with and without chunked prefill and full Llumnix on the A100 cluster.}
    \label{fig:llumnix_comparison}
\end{figure}

Figure~\ref{fig:llumnix_comparison} reports throughput, mean of request latency, token latency, and scheduling overhead. Astrolabe incurs a higher mean scheduling overhead than Llumnix at low load (23--29\,ms vs 5--7\,ms at QPS\,$\le$\,19) because it runs predictor queries and simulation before dispatch, but the overhead stays in the tens of milliseconds. Through QPS=19, both systems serve comparably; CP gives Astrolabe a small edge (5.5\,s vs 6.3\,s mean E2E at QPS=19), while the CP and no-CP Astrolabe curves remain within $\sim$10\% of each other. The high-load gap is therefore not explained by CP alone: Astrolabe without CP already outperforms Llumnix substantially. Llumnix collapses at QPS=20, where mean request latency jumps from 6.3\,s to 106\,s and mean token latency from 21\,ms to 1.15\,s. This is consistent with migration-induced PCIe/NIC contention: repeated KV-cache transfers can contend with normal serving traffic, allowing migration overhead to accumulate and dominate E2E latency. At QPS=36, the direct no-CP comparison already favors Astrolabe: Astrolabe without CP delivers 2.57$\times$ the throughput (18.7k vs 7.3k\,tok/s), 36$\times$ lower mean request latency (8.2\,s vs 300\,s), and 123$\times$ lower mean token latency (26\,ms vs 3.2\,s) than Llumnix. Enabling CP widens these to 2.58$\times$/40$\times$/135$\times$.
